\chapter{Mémoire de bord : décomposition paire--impaire et \(3\) réalisations hydrodynamique, de jauge et cosmologique}
\label{chap:memoria-tres-lecturas-rev6}
\index{mémoire de bord}\index{hydrodynamique}
\index{Yang--Mills}\index{cosmologie}

L'unité entre hydrodynamique, théorie de jauge et cosmologie ne s'obtient pas
en comparant trois équations une fois écrites. Elle apparaît auparavant, dans la
décomposition d'un même transport en une partie paire, qui conserve le coût
de l'aller et du retour, et une partie impaire, qui conserve son orientation.
L'objet commun est construit dans le système temporel dodécaphasique \(\mathcal R_{12}\) ; ses représentations
ultérieures ne s'identifient pas les unes aux autres.

\section{Des canaux angulaires au transport dodécaphasique}

Soient
\[
 q_\pm=\exp\!\left[-\frac{\pi}{180}(A\pm C^\ast)\right],
 \qquad
 S_{90}(q)=\frac{q^{90}}{1-q^{270}},
 \qquad
 S_{120}(q)=\frac{q^{120}}{1-q^{360}},
\]
et définissons
\[
 \kappa
 =12\bigl(S_{90}(q_-)-S_{90}(q_+)\bigr)
 -\bigl(S_{120}(q_-)-S_{120}(q_+)\bigr).
\]
Cette formule conserve la généalogie
\[
 \alpha_{\rm HMT}\longrightarrow A
 \longrightarrow C^\ast\longrightarrow q_\pm
 \longrightarrow\kappa.
\]
En particulier, \(\kappa\) n'est ni une viscosité, ni une tension, ni une masse. C'est
le coefficient adimensionnel de transfert que ces grandeurs pourront lire
ensuite.

Pour
\[
 \phi_m=(30m+20)^\circ,\qquad
 p_m=12\cos(3\phi_m)-\cos(4\phi_m),
 \qquad 0\leq m<12,
\]
sont posés
\[
 w_m=\frac{1+\kappa p_m}{12}.
\]
L'orthogonalité des caractères de
\(\mathbb Z/12\mathbb Z\) donne
\(\sum_m p_m=0\), donc \(\sum_mw_m=1\). En outre,
\(|p_m|\leq13\) ; pour la valeur construite de \(\kappa<1/13\), tous les
poids sont positifs. Sur
\(\mathcal H_{12}=\ell^2(\mathbb Z/12\mathbb Z)\), nous définissons
\[
 (U_{12}f)(n)=\sum_{m=0}^{11}w_mf(n-m).
\]

\begin{theorem}[Réserve dodécaphasique exacte]
\label{thm:reserva-dodecafasica-rev6}
Soient
\[
 P_{12}=\frac{U_{12}+U_{12}^{\ast}}2,\qquad
 J_{12}=\frac{U_{12}-U_{12}^{\ast}}2,\qquad
 L_{12}=I-P_{12}.
\]
Alors \(P_{12}\) est autoadjoint, \(J_{12}\) est
anti-autoadjoint et, pour tout \(f\perp\mathbf1\),
\[
 \boxed{\;
 \langle f,L_{12}f\rangle
 \geq g_{12}\|f\|_2^2,\qquad
 g_{12}=1-3\kappa>0.
 \;}
\]
La constante est optimale et est atteinte aux fréquences \(3\) et \(9\).
\end{theorem}

\begin{proof}
Dans la base de Fourier
\(e_k(n)=\exp(2\pi ikn/12)\), la convolution est diagonale :
\[
 U_{12}e_k=\lambda_k^{(U)}e_k,\qquad
 \lambda_k^{(U)}=\sum_{m=0}^{11}w_m e^{-2\pi ikm/12}.
\]
Le mode constant donne \(\lambda_0^{(U)}=1\). Les cosinus de fréquences trois et
quatre ne subsistent qu'en \(k=3,9\) et \(k=4,8\), respectivement ; les autres
modes s'annulent par orthogonalité. La plus grande valeur de
\(\operatorname{Re}\lambda_k^{(U)}\), avec \(k\neq0\), est \(3\kappa\). Par conséquent,
\[
 \langle f,L_{12}f\rangle
 =\sum_{k\neq0}(1-\operatorname{Re}\lambda_k^{(U)})|\widehat f(k)|^2
 \geq(1-3\kappa)\|f\|_2^2.
\]
\end{proof}

La décomposition possède une signification précise. \(P_{12}\) conserve ce
qu'une voie partage avec son inverse ; \(J_{12}\) conserve ce qui change de
signe lorsqu'on inverse l'orientation. Le premier peut fournir une forme
positive. Le second ne le peut pas : pour \(f\) réel,
\(\langle f,J_{12}f\rangle\) est purement imaginaire. En contrepartie,
\(J_{12}\) retient la circulation que la projection paire perdrait.

\section{Mémoire distribuée et générateur local}

La même dualité apparaît dans le noyau
\[
 K_q(m,n)=q^{|m-n|},\qquad 0<q<1.
\]
Sa matrice paraît non locale parce qu'elle corrèle des niveaux arbitrairement
séparés. Cependant, sur une chaîne finie \(0,\ldots,N\), son inverse
est tridiagonal :
\[
 K_q^{-1}
 =\frac1{1-q^2}
 \begin{pmatrix}
 1&-q&0&\cdots&0\\
 -q&1+q^2&-q&\ddots&\vdots\\
 0&-q&1+q^2&\ddots&0\\
 \vdots&\ddots&\ddots&1+q^2&-q\\
 0&\cdots&0&-q&1
 \end{pmatrix}.
\]

\begin{proposition}[Green non local, action locale]
\label{prop:green-local-memoria-rev6}
Pour \(0<q<1\), \(K_q\) est définie positive et la matrice précédente est son
inverse. Par conséquent, une mémoire distribuée exponentiellement est le
Green d'une action de premiers voisins.
\end{proposition}

\begin{proof}
Multiplier la matrice tridiagonale par \(K_q\) annule, dans chaque ligne intérieure,
les termes
\[
 -q\,q^{|m-1-n|}
 +(1+q^2)q^{|m-n|}
 -q\,q^{|m+1-n|}
\]
lorsque \(m\neq n\), et laisse \(1-q^2\) lorsque \(m=n\). Les lignes extrêmes
obéissent à la même identité avec un seul voisin. L'inverse est positif ;
par conséquent \(K_q\) l'est aussi.
\end{proof}

Ce résultat explique une opposition apparente qui traverse la mécanique
dimensionnelle : l'histoire est étendue parce qu'elle conserve les corrélations entre
profondeurs, tandis que le générateur qui la met à jour peut être local.
``Local'' et ``avec mémoire'' ne sont pas des alternatives.

\section{Représentations locales intérieure et réfléchie : circulation, holonomie, contraction et positivité de transfert}

Le chapitre du feuillet commun a démontré une seule fois les deux
représentations locales du bord. Dans la lecture intérieure, le secteur impair
porte la circulation et le secteur pair produit la contraction et la coercivité ; dans la
lecture réfléchie, le premier porte l'holonomie et le second la positivité de
transfert et l'écart fini. L'identité de Stokes et la coïncidence du
spectre non nul de \(d^\ast d\) et de \(dd^\ast\) sont une propriété de l'objet
commun, non deux preuves parallèles. Nous renvoyons à
\cref{thm:operador-comun-frontera-ns-ym,thm:contraccion-flujo-finito,%
thm:coercividad-gap-gauge-finito,thm:hoja-comun-frontera} ; nous ajoutons ici
uniquement la troisième lecture, la globale.

\section{Lecture globale : expansion et courbure}

La troisième lecture n'identifie pas l'univers à une cellule. Elle prend la même
séparation entre intérieur et bord comme loi globale d'évolution.
L'identité
\[
 1000=729+271
\]
rappelle que la croissance intérieure et la publication de bord complètent
une unité normalisée. Dans une géométrie homogène et isotrope, la partie paire
d'un transport peut être représentée par une pression effective
\[
 p_{\rm eff}=p-3\zeta H,
\]
qui est la forme tensorielle standard d'une viscosité de volume. L'équation
ne définit pas \(\zeta\) à partir de \(H\) ; elle déclare le type de la lecture cosmologique :
la mémoire symétrique est publiée comme résistance à l'expansion.

La géométrie de de Sitter montre, sans ambiguïté, ce qui demeure commun entre
trois feuilletages distincts. Pour \(H>0\),
\[
\begin{array}{c|c}
k&a_k(t)\\ \hline
+1&H^{-1}\cosh(Ht)\\
0&a_\ast e^{Ht}\\
-1&H^{-1}\sinh(Ht)
\end{array}
\]
dans leurs domaines respectifs.

\begin{proposition}[Une expansion, trois courbures spatiales]
\label{prop:tres-foliaciones-rev6}
Les trois fonctions précédentes satisfont
\[
 \frac{\dot a_k^2+k}{a_k^2}=H^2,\qquad
 \frac{\ddot a_k}{a_k}=H^2,
\]
et donc
\[
 \boxed{\quad
 R^{(4)}=12H^2,\qquad
 R^{(3)}=\frac{6k}{a_k^2}.
 \quad}
\]
\end{proposition}

\begin{proof}
Pour \(k=+1\), on utilise
\(\tanh^2+\operatorname{sech}^2=1\) ; pour \(k=-1\),
\(\coth^2-\operatorname{csch}^2=1\) ; la carte plane est immédiate.
Les formules de courbure de FLRW donnent alors les deux scalaires.
\end{proof}

Le signe \(k\) distingue la géométrie des feuillets spatiaux ; il ne change pas la
courbure de l'espace--temps de de Sitter. Cet exemple empêche de confondre
la lecture d'un feuillet avec celle de l'objet total, la même erreur que
l'architecture APP--TRIT--TPK évite dès le début.

\section{\(2\) représentations opératorielles exactes et \(1\) réalisation géométrique globale de la mémoire de bord}

\begin{theorem}[Une mémoire, deux représentations et une lecture globale]
\label{thm:una-memoria-tres-funciones}
La séparation paire--impaire du transport de bord admet deux
représentations opératorielles exactes et une lecture géométrique globale :
\[
\begin{array}{c|c|c|c}
\text{lecteur}&\text{secteur impair}&\text{secteur pair}&\text{statut}\\ \hline
\text{intérieur}&\text{circulation}&\text{contraction/coercivité}
&\text{représentation exacte}\\
\text{réfléchi}&\text{holonomie}&\text{noyau positif/écart fini}
&\text{représentation exacte}\\
\text{global}&\text{orientation du feuillet}&
                \text{réponse moyenne à l'expansion}
&\text{lecture FLRW typée}
\end{array}
\]
Les deux premières lignes partagent l'identité de Stokes et le spectre non
nul de \(d^\ast d\) et de \(dd^\ast\). La troisième conserve la distinction entre
courbure des feuillets et courbure globale, mais ne définit pas \(\zeta\) à partir de
l'opérateur dodécaphasique. Aucune ligne ne s'obtient en renommant les symboles d'une
autre.
\end{theorem}

\begin{proof}
La première et la deuxième lignes sont les
\cref{thm:contraccion-flujo-finito,thm:coercividad-gap-gauge-finito,thm:hoja-comun-frontera}.
La structure locale/non locale est prouvée dans la
\cref{prop:green-local-memoria-rev6}. La troisième ligne est réalisée par la
décomposition FLRW et la \cref{prop:tres-foliaciones-rev6} ; ces formules
prouvent la séparation géométrique, non un sélecteur autonome de la viscosité
de volume. Dans chaque cas, l'involution d'orientation fixe le canal impair et
la comparaison avec la voie inverse fixe le canal pair.
\end{proof}

La conclusion n'est pas qu'un fluide, une connexion et une cosmologie soient le
même système physique. Elle est plus précise : les trois sciences reçoivent
des représentations distinctes d'un unique système opératoriel de mémoire. Le flux la
lit vers l'intérieur, la théorie de jauge vers le bord et la cosmologie à
l'échelle globale.
